Cepheids are very bright variable stars whose mean absolute magnitudes are
functions of their pulsation periods. This allows astrophysicists to easily
determine their intrinsic luminosities from the variation in their observed,
apparent luminosities.

Below is a table with Cepheid data. $P_0$ is the pulsation period in days and
$\langle M_V\rangle$ is the mean absolute visual magnitude.

% The sixth Cepheid's name is printed with a corrupted glyph in the official
% paper (it appears as a CJK character); it is rendered here as alpha UMi.
\begin{center}
\begin{tabular}{ccc}
Cepheid & $P_0$ (days) & $\langle M_V\rangle$ \\
SU Cas & 1.95 & $-1.99$ \\
V1726 Cyg & 4.24 & $-3.04$ \\
SZ Tau & 4.48 & $-3.09$ \\
CV Mon & 5.38 & $-3.37$ \\
QZ Nor & 5.46 & $-3.32$ \\
$\alpha$ UMi & 5.75 & $-3.42$ \\
V367 Sct & 6.30 & $-3.58$ \\
U Sgr & 6.75 & $-3.64$ \\
DL Cas & 8.00 & $-3.80$ \\
S Nor & 9.75 & $-3.95$ \\
$\zeta$ Gem & 10.14 & $-4.10$ \\
X Cyg & 16.41 & $-4.69$ \\
WZ Sgr & 21.83 & $-5.06$ \\
SW Vel & 23.44 & $-5.09$ \\
SV Vul & 44.98 & $-6.04$ \\
\end{tabular}
\end{center}

\begin{enumerate}
\item Plot all Cepheids in a scatter diagram. $\log_{10}(P_0)$ should be the
  abscissa and $\langle M_V\rangle$ should be the ordinate.
\item Fit, using least squares, a straight line to the $\langle M_V\rangle$
  vs $\log_{10}(P_0)$ plot. This equation allows one to obtain the absolute
  magnitude from the pulsation period for any Cepheid.
\item Figures 1 and 2 show the light curves of two Cepheids. Use the
  available data to estimate the distances to each of these two Cepheids.
  Also estimate the uncertainty of the distance determination (exact
  formulae are not necessary).
\item Comparing the difference between the distances of the two stars with
  the typical size of a galaxy, would it be possible for these two stars to
  be in the same galaxy? Please mark ``YES'' (~) or ``NO'' (~).
\end{enumerate}

\paragraph{Appendix A: extinction coefficients.}
The magnitude outside the atmosphere is given by:
\[ M = m - A\cdot X + B \]
where $A$ is the extinction coefficient and $B$ is the zero point coefficient
for the night, $X$ is the airmass of the observation and $m$ is the
instrumental magnitude (that is, the magnitude obtained directly from the
image).

\paragraph{Appendix B: calibrated differential magnitude.}
If the standard star is in the same frame as the object, the calibrated
magnitude of the object can be obtained by:
\[ M_{\mathrm{ast}} = m_{\mathrm{ast}} - m_{\mathrm{star}}
   + M_{\mathrm{star}} \]

\paragraph{Appendix C: least squares estimation of angular coefficients.}
Given straight lines described by
\[ y_i = \alpha + \beta x_i,\ i = 1..n \]
the least squares estimator of $\beta$ is given by
\[ \beta = \frac{\sum_i^n x_i y_i - \frac{1}{n}\sum_i^n x_i \sum_i^n y_i}
   {\sum_i^n x_i^2 - \frac{1}{n}\left(\sum_i^n x_i\right)^2} \]
